\section{Post-processing}\label{sec:postproc}

A completed run leaves behind the log-scale history $\ln U_k, \ln L_k$ and the
clock history $\ln \tau_k$. From these we recover the power-law parameters
$\alpha_X, \beta_X$ of the ansatz \eqref{eq:fw-powerlaw}. This is the offline,
global counterpart of \textsc{EstimateBeta} (\S\ref{ssec:alg-beta-estimation}),
which fits $\beta_X$ one step at a time; here the whole history is used at
once, so the fit also serves as a check on the run.

Write
\begin{equation}\label{eq:pp-logs}
  \mathcal{L}_k = \ln L_k, \qquad \mathcal{U}_k = \ln U_k,
\end{equation}
together with the per-step log-ratios
\begin{equation}\label{eq:pp-log-ratios}
  \ell_k = \ln L_{k+1} - \ln L_k, \qquad \lambda_k = \ln U_{k+1} - \ln U_k,
\end{equation}
which for a Lie--Trotter step equal $\ln \tilde{L}_{k+1}$ and
$\ln \tilde{U}_{k+1}$ of \eqref{eq:fw-ratios}, and for a Strang step combine
both half-step ratios. The stored scales are their cumulative sums,
\begin{equation}\label{eq:pp-cumsum}
  \mathcal{L}_k = \mathcal{L}_0 + \sum_{i=1}^{k} \ell_i,
  \qquad
  \mathcal{U}_k = \mathcal{U}_0 + \sum_{i=1}^{k} \lambda_i,
\end{equation}
which grow only linearly in $k$; unlike $L_k$ and $U_k$ themselves,
$\mathcal{L}_k$ and $\mathcal{U}_k$ never overflow.

\subsection{The scaling rates under the ansatz}\label{ssec:pp-rates}

Differentiating the ansatz \eqref{eq:fw-powerlaw} with respect to $t$, and
using the clock \eqref{eq:fw-clock} to convert from $t^*$, the logarithmic
rates \eqref{eq:fw-rates} are
\begin{equation}\label{eq:pp-rate-chain}
  \sigma_X = \frac{1}{X}\frac{\mathrm{d} X}{\mathrm{d} t^*}\frac{\mathrm{d} t^*}{\mathrm{d} t}
  = \frac{\beta_X\, \tau}{t^* + \alpha_X},
  \qquad X \in \{U, L\}.
\end{equation}
Inverting the ansatz to eliminate the unknown offset,
\begin{equation}\label{eq:pp-tstar-offset}
  t^* + \alpha_X = X^{1/\beta_X},
\end{equation}
and substituting into \eqref{eq:pp-rate-chain} gives the rates in terms of the
scales alone,
\begin{equation}\label{eq:pp-rate}
  \sigma_X = \beta_X\, \tau\, X^{-1/\beta_X}.
\end{equation}
Taking logarithms,
\begin{equation}\label{eq:pp-rate-log}
  \ln \sigma_{X,k} = \ln \beta_X + \ln \tau_k - \frac{\mathcal{X}_k}{\beta_X},
  \qquad \mathcal{X}_k = \ln X_k .
\end{equation}
Finally, the rates themselves are estimated directly from the log-ratios
already in hand, without any finite difference of $L_k$ or $U_k$ themselves,
\begin{equation}\label{eq:pp-rate-fd}
  \sigma_{L,k} \approx \frac{\ell_k}{\Delta t},
  \qquad
  \sigma_{U,k} \approx \frac{\lambda_k}{\Delta t}.
\end{equation}

\subsection{Recovering \texorpdfstring{$\alpha_X$}{alpha\_X} and \texorpdfstring{$\beta_X$}{beta\_X}}\label{ssec:pp-fit}

Equation \eqref{eq:pp-rate-log} is linear in $\mathcal{X}_k$ once the clock is
divided out. Regressing
\begin{equation}\label{eq:pp-regression}
  \ln \sigma_{X,k} - \ln \tau_k = \ln \beta_X - \frac{1}{\beta_X}\, \mathcal{X}_k
\end{equation}
against $\mathcal{X}_k$ therefore yields $\beta_X$ twice over: once as
$e^{\text{intercept}}$ and once as $-1/\text{slope}$. Agreement between the two
is a check that the run is in fact following a power law over the fitted range;
disagreement indicates that it is not, or that the range includes the initial
transient before the frame has locked on.

With $\beta_X$ in hand, the offset follows from \eqref{eq:pp-tstar-offset},
\begin{equation}\label{eq:pp-alpha}
  \alpha_X = X_k^{1/\beta_X} - t_k^*,
\end{equation}
where the physical time is accumulated from the per-step increments of
\eqref{eq:fw-dtstar},
\begin{equation}\label{eq:pp-tstar-cumsum}
  t_k^* = t_0^* + \sum_{i=0}^{k-1} \Delta t_i^* .
\end{equation}
In exact arithmetic \eqref{eq:pp-alpha} is independent of $k$; in practice it
is averaged over the fitted range, and its scatter is a second check on the
ansatz.

For the 1D diffusion equation the clock is $\tau = L^2/D$, and
\eqref{eq:pp-rate-log} specializes to
\begin{equation}\label{eq:pp-1d-sigma-L}
  \ln \sigma_{L,k} = \ln \frac{\beta_L}{D} + \left( 2 - \frac{1}{\beta_L} \right) \mathcal{L}_k,
\end{equation}
\begin{equation}\label{eq:pp-1d-sigma-U}
  \ln \sigma_{U,k} = \ln \frac{\beta_U}{D} + 2\, \mathcal{L}_k - \frac{\mathcal{U}_k}{\beta_U},
\end{equation}
so that for $X = L$ the regression \eqref{eq:pp-regression} collapses onto a
single variable, while for $X = U$ it is a two-variable fit in
$\mathcal{L}_k$ and $\mathcal{U}_k$.
